
\chapter{Correlation, the NumPy library}
\label{ch:correlation}

\section{``Are these things correlated?''}

A very common exploratory operation is to ask yourself whether two variables
in a data set are related to each other in some way. After looking at
variables individually (via things like min/median/mean/max, histograms, bar
charts) this is often the next thing you do. So how can we tell whether there
is a genuine relationship between variables, versus just a spurious
pseudo-connection?

\subsection{Statistical significance}

This question is at the heart of \textbf{statistical significance}, one of the
most important concepts in data analysis. In general, when you ask whether your
data shows evidence of some pattern, you need to determine whether that pattern
is pronounced enough that it is ``statistically significant'' (or ``stat sig,''
for short.) This essentially means that the pattern (in this chapter,
the correlation) is so strong that it's very unlikely for it to have been a
fluke of your particular data set.

It's clearer with an example. Suppose you wanted to know: ``do cats or dogs
have longer \textit{names}?'' To answer this, you go collect some data from
your friends who have pets. The cat names you collect are ``\textsf{Felix},''
``\textsf{Garfield},'' ``\textsf{Smudge},'' ``\textsf{Pillow},'' and a few
others. The dog names you collect are ``\textsf{Coco},'' ``\textsf{Fido},''
``\textsf{Mrs. Muggles},'' ``\textsf{Rex},'' and a few others. Ever the Data
Scientist, you take the average. You discover that in your data set, the mean
number of letters in a cat name is 6.8, and the mean number of letters in a dog
name is only 6.3. So cats \textit{in your sample} have longer names. But the
question that really matters is: ``\textit{generally speaking, on a consistent
and reliable basis, will cats have longer names on average than dogs will?}''

This is a question about more than just your friends' pets. It's a question
about \textit{all} pets, of which you have only sampled a tiny fraction. But
normally, that's the question you actually care about.

Now there are three different ways to answer this question, two of them
\textbf{wrong}, and one of them \textbf{right}:

\begin{itemize}
\itemsep.1em
\item \textbf{Wrong way \#1}: Look at your numbers and say, ``wow, the cats I
sampled had on average 6.8 letters as opposed to just 6.3. That sure seems like
a big difference to me! I declare that cats have longer names in general.''
\item \textbf{Wrong way \#2}: Look at your numbers and say, ``meh, my cat names
only beat my dog names by 0.5. That sure doesn't seem like a very big
difference to me! I declare that overall, cats and dogs are about equal.''
\item \textbf{Right way}: Run the appropriate statistical test, and go with its
results.
\end{itemize}

The majority of humans I've ever met choose one of the first two ways.
(Normally they choose whichever way supports the conclusion they'd like to
draw.) But we are Data Scientists, and we will always choose the third way. 

All we need to know is what the appropriate statistical test \textit{is} for a
given situation, and how to interpret its results. For detecting correlations
between variables, that in turn depends on what \textit{types} of data you're
dealing with (categorical or numerical, for instance). There are three cases:
two categorical variables, one categorical and one numeric variable, and two
numeric variables. Let's look at each in turn.


\subsection{Two categorical variables}

As explained in section~\ref{contingencyTable} (p.~\pageref{contingencyTable}),
the first step in analyzing two categorical variables is creating a contingency
table. Then, the next question is usually, ``okay, the proportions of the
values of one variable are different across the other variable...but is this
statistically significant?'' A quick and dirty way to tell is to run a
\textbf{$\chi ^2$ test} (sometimes written ``chi-squared,'' but pronounced
``kai squared''), as shown below.

\label{chisq}

\fbox{\parbox{\textwidth}{

\textbf{*WARNING!} When you run a $\chi ^2$ test on a contingency table, as
shown below, \textit{it must be a \textbf{margin-less} contingency table!!} In
other words, if you used any of the techniques from section~\ref{marginsBegin}
(pp.\pageref{marginsBegin}-\pageref{marginsEnd}) in order to see
row/column/overall totals, you must \textbf{\textit{not}} use that table with
the extra ``\texttt{All}'' columns as the argument to the
\texttt{chi2\_contingency()} function described below. Doing so will give you
\textit{wrong} and misleading answers, with no obvious indication that they are
wrong. Always give \texttt{chi2\_contingency()} the original contingency table,
\textit{without} margins.

}}

\bigskip
\medskip
Okay. Warnings aside, here's the code to perform the $\chi ^2$ test:

\begin{Verbatim}[fontsize=\small,samepage=true,frame=single]
>>> import scipy.stats
>>> scipy.stats.chi2_contingency(my_contingency_table)
\end{Verbatim}

\begin{samepage}
For example, using our Star Wars characters DataFrame:

\begin{Verbatim}[fontsize=\small,samepage=true,frame=single]
>>> side_saber = pd.crosstab(characters.side, columns=characters.saber)
>>> side_saber
    saber  blue  green  red
    side                   
    Dark      0      0    2
    Light     2      2    0

>>> scipy.stats.chi2_contingency(side_saber)
    (6.0, 0.049787068367863951, 2, 
            array([[ 0.66666667,  0.66666667,  0.66666667],
            [ 1.33333333,  1.33333333,  1.33333333]]))
\end{Verbatim}
\end{samepage}

Whew -- that's some strange stuff that \texttt{chi2\_contingency()} printed
out. It's a list\footnote{Technically, it's a Python ``\textbf{tuple}'' rather
than a list, as indicated by bananas (parentheses) instead of boxies (square
brackets) around the outside of it. But to analyze it, you can treat it as a
list. In particular, you can use ``\texttt{len()}'', ``\texttt{[0]}'',
``\texttt{[1]}'', ``\texttt{[2]}'',
\textit{etc.}~to count and extract the different elements.} of four elements.
Briefly, here's what these elements mean:

\begin{compactenum}
\item[0.] The raw value of the actual ``$\chi ^2$ statistic.'' \textit{Ignore}
this.
\item[1.] The test's ``\textbf{p-value},'' which is by far the most important of
the four elements. As with many statistical tests that return a ``p-value,'' if
it's less than \textbf{.05}, you can be confident that the two variables
\textit{are} correlated in some way. (In the above example, they are...barely.)
\item[2.] The test's number of ``degrees of freedom.'' \textit{Ignore} this.
\item[3.] The ``expected frequencies,'' which you also usually \textit{ignore}, but
for the curious it is sort of interesting. Notice that that fourth element
(element \#3 of the list) is a thing called an \texttt{array}, which is
actually a two-dimensional \textbf{NumPy nd-array} (see section~\ref{NumPy},
below). If you stare at it, you'll realize that it's 6 numbers arranged in a
box: three on the top row (0.667, 0.667, and 0.667) and three on the bottom
(1.33, 1.33, and 1.33). This tells you how many of each of the
lightsaber color / dark-side-light-side combinations we would have expected on average
\textit{if lightsaber color had no correlation with side.} Here, it's saying,
``if the color lightsaber you choose really has nothing to do with what side of
the Force you use, and you had six total people (two blue, two green, and two
red) then
we'd expect there to be about $\frac{2}{3}$rds of a blue-sabered Light-side
person, $\frac{2}{3}$rds of a green-sabered Light-side
person,
about $\frac{2}{3}$rds of a red-sabered Light-side person, $1\frac{1}{3}$rd
blue-sabered Dark-side people, 
$1\frac{1}{3}$rd green-sabered Dark-side people, and
$1\frac{1}{3}$rd red-sabered Dark-side people.'' The $\chi ^2$ test is basically
looking to see how far off the \textit{actual} numbers of the various
saber/side counts are from this baseline expectation. If they're off by a
lot, the p-value (in element \#1) will be less than .05, telling us that
lightsaber color \textit{does} have something to do with what side of the Force
you use.
\end{compactenum}

Okay, that was a lot. The short version of the above: ignore everything except
element \texttt{[1]} of the \texttt{chi2\_contingency} output. If that number
is less than .05, interpret it as ``these two categories variables \textit{are}
in fact correlated,'' and if it's higher, interpret it as ``they don't look to
be.''

\subsection{One categorical and one numeric variable}

My favorite tool for one categorical and one numeric variable is the
\textbf{notched boxplot} from section~\ref{notched} on p.~\pageref{notched}.
(You can also do a t-test if the notches look funny; see
\texttt{scipy.stats.ttest\_ind} and \texttt{scipy.stats.ttest\_rel} and read
carefully.)


\subsection{Two numeric variables}

Besides visually inspecting scatterplots to spot correlations, the Pearson's
correlation coefficient can give you a quick-and-dirty answer:
\label{pearsonsR}

\begin{Verbatim}[fontsize=\small,samepage=true,frame=single]
>>> import scipy.stats.stats
>>> scipy.stats.stats.pearsonr(my_var1, my_var2)
\end{Verbatim}

For instance:

\begin{Verbatim}[fontsize=\small,samepage=true,frame=single]
>>> scipy.stats.stats.pearsonr(men_ht,men_sat)
    (-0.080877008297474018, 0.29443331215713264)
\end{Verbatim}

The second number is the p-value: if it's greater than about .05, then these
two variables are likely \textit{not} correlated.


\section{NumPy}
\label{NumPy}

At the near-bottom of the Python data analysis ecosystem is the NumPy
library. SciPy, Pandas, scikit-learn, and many other things are built on top
of it.

Its wheelhouse is the \textit{ndarray}, a highly-optimized and efficient
multidimensional array datatype. It's blazing fast, and supports advanced
vectorized and broadcasting operations as we'll see. A lot of this boils down
to \textit{avoiding writing expensive loops in Python} and letting the faster,
native, underlying language (C) do bulky computations in a fraction of the
time.

NumPy also has a lot of other really useful stuff in it, including linear
algebra routines and sophisticated random number generation. In practice, you
often use some higher-level library to access these things, but under the hood
they in turn use NumPy.

\subsection{ndarrays}

The standard NumPy import is:

\begin{Verbatim}[fontsize=\small,samepage=true,frame=single]
import numpy as np
\end{Verbatim}

after which you can make \textit{ndarray}s in a variety of ways:

\begin{Verbatim}[fontsize=\footnotesize,samepage=true,frame=single]
a = np.array(['larry','curly','moe'])    # 1-d array ("vector")
b = np.array([[1,2,3],[8,9,10])          # 2-d array ("matrix")
c = np.zeros(25)                         # 1-d array of all zeros
d = np.zeros((9,7))                      # 2-d array of all zeros
                                         # See also ".ones()" and ".empty()"
e = np.arange(5,855)                     # Like .range(), but for ndarrays

\end{Verbatim}

\subsubsection{dtypes}

Every \textit{ndarray} has a \texttt{dtype}, or ``data type,'' which you can
access using the ``\texttt{.dtype}'' accessor. You can also call
``\texttt{.astype(...)}'' to perform data conversion. The types we'll care most
about are \texttt{np.int32} and \texttt{np.float64} for integers and
floating-point (``decimal'') numbers, respectively.

\subsubsection{vectorization / broadcasting}

Similar to how you've seen Pandas work, NumPy provides vectorizing and
broadcasting operations so you can avoid writing and executing loops. For
example, if you want to add the elements of two arrays together, instead of
doing something like this:

\begin{Verbatim}[fontsize=\small,samepage=true,frame=single]
for i in range(len(m)):
    q[i] = m[i] + n[i]
\end{Verbatim}

you just write this:

\begin{Verbatim}[fontsize=\small,samepage=true,frame=single]
q = m + n
\end{Verbatim}

Not only is this nicer syntactically, and less typing (and opportunity for
error), but it's \textit{much, much faster} for large data sets. It cannot be
overemphasized that you want to avoid loops (and especially nested loops) for
any serious calculations if you can help it at all. Sometimes this requires
rewiring your brain a bit to anticipate where you can take advantage of
vectorized operations.

Play around with this idea using operators like ``\texttt{**}'' and
``\texttt{<=}''.

In addition to operators, NumPy also provides a number of functions (called
``\textit{ufuncs}'' or ``universal functions'') that work in a similar way.
Instead of:

\begin{Verbatim}[fontsize=\small,samepage=true,frame=single]
for i in range(len(f)):
    g[i] = math.sqrt(f[i])
\end{Verbatim}

you just write this:

\begin{Verbatim}[fontsize=\small,samepage=true,frame=single]
g = np.sqrt(f)
\end{Verbatim}

Play around with ufuncs like ``\texttt{np.sin()}'', and
``\texttt{np.power()}''.

There are also many other standard functions which take an \textit{ndarray}
and return a single result, like ``\texttt{np.max()}'', ``\texttt{np.any()}'',
``\texttt{np.cumsum()}'', and many, many others. \textit{Before you ever write
a loop, take a moment to consider whether it might be a good idea to browse
the NumPy documentation and avoid it.}

\subsubsection{Reshaping}

NumPy \textit{ndarray}s have a \texttt{.reshape()} function that lets you
coerce a multi-dimensional array to different dimensions:

\begin{Verbatim}[fontsize=\small,samepage=true,frame=single]
thing = np.array([[1,2,3],[4,5,6],[7,8,9],[10,11,12]])
print(thing)
[[ 1  2  3]
 [ 4  5  6]
 [ 7  8  9]
 [10 11 12]]

thing = thing.reshape(2,6)
print(thing)
[[ 1  2  3  4  5  6]
 [ 7  8  9 10 11 12]]

thing = thing.reshape(3,4)
print(thing)
[[ 1  2  3  4]
 [ 5  6  7  8]
 [ 9 10 11 12]]
\end{Verbatim}

You can also perform a matrix \textbf{transpose} operation using \texttt{.T}
(no bananas, note):

\begin{Verbatim}[fontsize=\small,samepage=true,frame=single]
print(thing.T)
[[ 1  5  9]
 [ 2  6 10]
 [ 3  7 11]
 [ 4  8 12]]
\end{Verbatim}

\subsubsection{Slices: warning}

You can take ``slices'' of an \textit{ndarray} just like you can a normal
Python list, using the colon syntax. Beware, though: \textit{in NumPy, these
are \textbf{views} on the original list, not copies!} In order to make a copy,
you must explicitly call ``\texttt{.copy()}'' on the array.

To illustrate the possible weirdness:

\begin{Verbatim}[fontsize=\footnotesize,samepage=true,frame=single]
x = np.array([4,9,8,3,2])
y = x[2:4]
print(y)
[8 3]

y[1] = 999999999
print(x)
[        4         9         8 999999999         2]

\end{Verbatim}

\subsection{Random number generation}

One thing we'll use surprisingly often in Data Science is random number
generation. NumPy provides a super good library for this, all of which is in
the namespace \texttt{np.random}. Just about every distribution you can think
of is present, including:

\begin{Verbatim}[fontsize=\small,samepage=true,frame=single]
print(np.random.normal(10,3,5))
[ 13.14637747   8.0170962   14.69793015   5.28517005   7.71561529]

print(np.random.normal(10,30,5))
[ -7.17146      9.32914316   7.63719265  11.72634167  23.12532128]

print(np.random.uniform(6,10,5))
[ 7.3502029   8.82037522  7.89682321  7.75582494  7.98879983]

print(np.random.exponential(5,5))
[ 2.41409743  0.59398243  2.12184356  2.04372147  1.56823102]

print(np.random.exponential(50,5))
[ 47.4720236    6.82838706   5.14356732  60.47055961   0.47463824]
\end{Verbatim}

In all these examples, the last argument to the function is the desired
\textit{size} (number of elements) of the resulting array.

